\scanpage{76}
$\unsure{\text{canon}^{t}.|\text{canon}^{t}|\text{canonit.}}$ isom à un polygône combinatoire-type
$\Pi_n$ $(3\leq n\leq+\infty)$. D'ailleurs
tous ces $\Pi_n$ sont \inserted{\unsure{canoniqut|canoniques|canoniquement}} quotients du
même $\Pi_\infty$ (si $n\mid n'$, $\Pi_n$ est quotient
de $\Pi_\infty$), $r_n$ \inserted{le rep. de $\Pi_n$} étant l'image du repère $r_\infty$ de
$\Pi_\infty$, donc la cat. (discrète) des repr. affines
régulières de $\Pi_n$ se plonge (comme
sous-catégorie pleine) dans celle des repr
régulières de $\Pi_\infty$. Donc le plus
pratique sera pour \unsure{ns|nous} de faire l'étude
en trois étapes :

a) Étudier les réal. géom. rég. de $\Pi_\infty$
sur un corps (plus gén anneau commutatif $k$)
\deleted{\unsure{commutatif|un invariant}} quelconque. On verra entre autres
qu'ils sont \unsure{catrctérisés|caractérisés|caractérisées}, à isom près, par un
l'invariant numérique $\alpha\in k$, qui peut
être choisi arbitrairement.

b) Expliciter les conditions sur $\alpha$
pour que la réal. géom envisagée de
$\Pi_\infty$ se factorise par $\Pi_n$, et plus
% archive sheet number: 75
